\section{Uniformity and commutator estimates}
\label{app:comparison}

The proof distinguishes ordinary norm control, control at every matrix
size, and implementation on the original Hilbert space. We compare the
uniformity required for these steps with the commutator estimate in the
similarity companion, and then record how the small-unitary conclusion
controls both matrix gaps and commutants.

\subsection{Uniformity over factors and representations}

Fix an abstract type $\II_1$ factor $M$ with separable predual.
A stability statement uniform over its representations would assign,
for each $\eps>0$, a number $\delta_M(\eps)>0$ that works for every
faithful normal unital representation $\rho:M\to\B(H)$ and every
unital von Neumann algebra $N\subseteq\B(H)$: if
$d(\rho(M),N)<\delta_M(\eps)$, then a unitary within $\eps$ of the
identity conjugates $\rho(M)$ onto $N$. Even with this representation
uniformity, the tolerance may depend on $M$.

The quantifier order for universal stability is stronger:
the tolerance must be chosen before the factor as well as before its
representation. Positivity of each $\delta_M(\eps)$ does not imply
\[
 \inf_M\delta_M(\eps)>0.
\]
Consequently, a representation-transfer argument with absolute constants
preserves an already universal tolerance but does not by itself remove
factor dependence from its input.

In \Cref{thm:finite}, uniformity is obtained before the concluding
reductions. The common-position theorem has absolute constants, the
boundary argument in \Cref{prop:finite-standard-isomorphism} gives an
isomorphism with an absolute ordinary-norm bound, and
\Cref{thm:amplification} supplies an absolute amplification constant.
Their combination gives the single modulus $\eta_{\rm fin}$ for all
separably acting finite factors. Sections~\ref{sec:cores}
and~\ref{sec:reductions} propagate this uniform control through finite
corners, measurable assembly, and separable restrictions.

\subsection{Two commutator arguments}

The amplification theorem needed here concerns selfadjoint implementers
of commutators on separably acting type $\II_1$ factors. The similarity
companion proves a broader estimate: for every unital von Neumann algebra
$P\subseteq\B(K)$, every $Y\in\B(K)$, every integer $n\ge1$, and every
$X\in M_n(P)$,
\[
 \norm{[Y^{(n)},X]}
 \le C\norm{\ad_Y|_P}\norm X,
\]
where $Y^{(n)}$ is the diagonal amplification on $K^n$ and $C$ is
absolute, with no separability assumption
\cite[Theorem~1.2]{OpenAI2026Similarity}.

Both arguments use cyclic representation estimates and tracial free
products to obtain bounds independent of matrix size. In the present
paper, the column estimate of \Cref{lem:column-commutator} controls
randomized matrix coefficients. The free-word estimates bound their
commutators, while the averaged products converge weakly to a positive
scalar multiple of the identity in \eqref{eq:weak-scalar-limit}.
Together these facts yield \Cref{thm:amplification}.

The companion's central intermediate result is a corner estimate.
It works in a tracial free product $\mathcal L=D*A_0$, where $D$ is a
hyperfinite $\II_1$ factor generated by increasing full matrix algebras
and $A_0$ is a finite von Neumann algebra with a faithful normal tracial
state. Equip $\mathcal L$ with the free-product trace.
Writing $L_x$ for left multiplication by $x$ on $L^2(\mathcal L)$,
suppose that $G\in\B(L^2(\mathcal L))$ commutes with $L_D$ and that,
for some $e\ge0$,
\[
 \norm{[G,L_x]}\le e\norm x\qquad(x\in p\mathcal Lp),
\]
where $p\in D$ is a projection of trace $1/m$ for an integer $m\ge2$. For a
trace-zero selfadjoint unitary $s\in A_0$, the companion bounds the
compression of $[G,L_s]$ to $L^2(A_0)$ by $C_0e$, with $C_0$
independent of $m$ \cite[Theorem~3.1]{OpenAI2026Similarity}.
Its proof analyzes invariant tensors of reduced words, approximates
operators by right multipliers on small right supports, and obtains the
trace-independent bound from a Hankel compression
\cite[Sections~3--5]{OpenAI2026Similarity}. Using Popa's independence theorem,
the companion transfers this estimate to finite factors; passage through a
countable amplification treats their arbitrary separable normal
representations \cite[Section~6 and Proposition~7.1]{OpenAI2026Similarity}.

The finite stability proof here uses \Cref{thm:amplification} directly.
For the near-identity isomorphism $\theta:M\to N$ constructed there,
the switch operator on $H\oplus H$ has commutator with
$x\mapsto x\oplus\theta(x)$ whose off-diagonal entries are
$x-\theta(x)$ and its negative. Applying the amplification theorem to
this paired representation converts the ordinary bound on $\theta-\id$ into a
cb bound. \Cref{lem:paired} then supplies the unitary on the original
$H$. This identifies the precise role of commutator control in the
stability argument.

\subsection{Matrix gaps and commutants at the conclusion}

\begin{corollary}\label{cor:uniform-matrix-commutant}
For every $\eps>0$ there is $\delta>0$, independent of the algebras,
their representations, and their Hilbert spaces, such that unital von
Neumann algebras $M,N\subseteq\B(H)$ with the same identity and
$d(M,N)<\delta$ satisfy
\[
 \dcb(M,N)<\eps,\qquad d(M',N')<\eps.
\]
\end{corollary}

\begin{proof}
Choose $\delta$ from \Cref{thm:main} with unitary displacement
$\eps/2$. The resulting unitary satisfies $uMu^*=N$ and
$\norm{u-I}<\eps/2$. Conjugation by its diagonal amplification
identifies the unit balls of $M_n(M)$ and $M_n(N)$ and moves each
contraction by at most $2\norm{u-I}$, independently of $n$.
The same estimate applies to the commutants because $uM'u^*=N'$.
\end{proof}

Since $d(M,N)\le\dcb(M,N)$, this corollary shows that the ordinary
and completely bounded distances induce the same uniform structure on
the unital von Neumann algebras in each $\B(H)$, with tolerances uniform
over $H$.
It also gives uniform continuity of the commutant operation for the
ordinary distance. These controls follow at the conclusion from one
small implementing unitary; the separate estimates in the proof are
what make construction of that unitary possible.
